\section{Short-Circuit Models}

\subsection{Problem Form and Runtime Scope}

This chapter covers the short-circuit analysis layer built on a projected
canonical AC network.  The core computation is a fault-network Thevenin solve
for one or more faulted buses, with both simplified and detailed IEC-style
reporting routes exposed by the current API.

\subsection{Code Map}

The current implementation is concentrated in:
\begin{itemize}
  \item \texttt{src/analysis/short\_circuit.cpp} for batch and single-bus fault execution,
  \item canonical network projection and admittance assembly reused from the main power-flow preprocessing stack,
  \item the sparse linear-solver abstraction in the engine layer for repeated fault-network solves.
\end{itemize}

\subsection{Pseudocode Map}

Section~\ref{sec:algorithms} provides the code-aligned workflow summary.
Algorithm A12 is the direct pseudocode reference for this chapter.

\subsection{Current Status}

The detailed short-circuit route is now executable rather than placeholder-only.
The remaining gap is mostly in result-contract cleanup and richer external-grid
reporting, not in the existence of a working solver path.

\subsection{Implemented Solver Paths}

The short-circuit module currently exposes two public result families:
\begin{itemize}
  \item \texttt{compute\_short\_circuit}: batch computation over all buses or a user-selected subset,
  \item \texttt{compute\_fault\_at\_bus}: single-bus query using the same underlying fault-admittance model.
\end{itemize}

In addition, the module now exposes an executable detailed IEC-style route under
\texttt{SCDetailedOptions} / \texttt{SCDetailedResult} through
\texttt{run\_short\_circuit\_detailed} and
\texttt{run\_short\_circuit\_detailed\_batch}.  That detailed path computes
\(I_{k}^{\prime\prime}\), \(i_p\), \(I_b\), and \(I_k\) for all four supported
fault types.  The public struct still carries legacy scaffold fields
\texttt{placeholders\_only} and \texttt{pending\_functions}, but the current
implementation sets them to \texttt{false} and empty and fills the detailed
per-bus current table.

\subsection{Fault-Network Construction}

The short-circuit entry point first projects the rich information model to the
canonical AC network, then constructs a fault admittance matrix from the
following components:
\begin{align}
Y_{fault} = Y_{branch} + Y_{gen}^{\prime\prime} + Y_{motor} + Y_{load,motor} + Y_{ext}.
\end{align}

\paragraph{Branch contribution.}
All in-service \texttt{ACBranch} objects, including projected transformers and
switches, contribute their positive-sequence admittance to \(Y_{branch}\).  Tap
ratios and phase shifts are handled through the complex off-nominal transformer
model.

\paragraph{Generator subtransient contribution.}
Each in-service generator contributes a shunt term approximately equal to
\begin{align}
Y_{gen}^{\prime\prime} \approx \frac{1}{j x_d^{\prime\prime}},
\end{align}
using \texttt{Generator::xdpp\_pu} when present and a fallback value from
\texttt{SCOptions::default\_xdpp} otherwise.

\paragraph{Motor contribution.}
The solver adds subtransient motor admittances from two sources:
\begin{itemize}
  \item explicit \texttt{AsynchronousMotor} objects in \texttt{ACSystem::motors},
  \item the motor fraction embedded in \texttt{Load} objects through \texttt{motor\_percent}, \texttt{x\_sub\_pu}, and \texttt{r\_sc\_pu}.
\end{itemize}

\paragraph{External grid contribution.}
External-grid data is used as a source equivalent in the fault network using
available impedance or short-circuit-strength fields.

\subsection{Z-Bus Evaluation}

Once \(Y_{fault}\) is assembled, the solver factorizes it and obtains the
Thevenin impedance at a faulted bus \(k\) from the \(k\)-th diagonal entry of the
inverse:
\begin{align}
Z_{kk} = \left(Y_{fault}^{-1}\right)_{kk}.
\end{align}
Rather than building the full inverse explicitly, the implementation solves
\begin{align}
Y_{fault} z = e_k
\end{align}
for each faulted bus and reads the \(k\)-th component of the solution vector.

\subsection{Current Implemented Fault Equations}

The current simplified solver assumes a prefault voltage of
\begin{align}
V_0 = c\angle 0, \qquad c = \texttt{SCOptions::c\_factor}.
\end{align}

\paragraph{Three-phase fault.}
\begin{align}
I_f^{3\phi} = \frac{V_0}{Z_{kk}}.
\end{align}

\paragraph{Single-line-to-ground, line-to-line, and double-line-to-ground faults.}
The current code uses simplified sequence-network approximations when explicit
negative- and zero-sequence detail is unavailable.  In effect,
\begin{itemize}
  \item the negative-sequence network is approximated by the positive-sequence network,
  \item the zero-sequence network is approximated or simplified rather than fully assembled from every device model.
\end{itemize}
Therefore, the unsymmetrical-fault results are best interpreted as engineering
approximations rather than a complete IEC 60909 sequence implementation.

\subsection{Reported Quantities}

For each bus, the primary result object stores:
\begin{align}
S_k^{sc} &= \frac{S_{base}|V_0|^2}{|Z_{kk}|}, \\
I_{k}^{\prime\prime} &= |I_f| I_{base},
\end{align}
where the current is reported in kA via the bus voltage base.

\subsection{Detailed IEC Helper Functions}

The detailed SC header also defines helper routines for:
\begin{itemize}
  \item voltage-factor selection,
  \item \(\kappa\)-factor estimation from \(R/X\),
  \item generator and transformer correction factors,
  \item zero-sequence correction factors,
  \item motor-impedance reconstruction.
\end{itemize}
These helpers are now used by the detailed solver path directly.  The remaining
gap is no longer ``can the detailed route run?'' but rather whether the public
result contract exposes every contribution channel explicitly enough for
downstream consumers.

\subsection{Practical Interpretation}

For the current code base, the short-circuit module should be interpreted as:
\begin{itemize}
  \item \implemented{} for positive-sequence Z-bus short-circuit studies with generator, motor, and external-grid contributions,
  \item \implemented{} for approximate unsymmetrical-fault calculations,
  \item \implemented{} for detailed IEC-style current calculations \(\{I_{k}^{\prime\prime}, i_p, I_b, I_k\}\) across 3\(\phi\), SLG, LL, and DLG faults,
  \item \planned{} for a cleaner detailed public contract with explicit external-grid contribution fields and removal of the leftover scaffold markers.
\end{itemize}
